\subsection{The arithmetic boundary norm and the compact orientation class}

\begin{statement}{D.55}{Proved}{An exact morphism from the arithmetic boundary obstruction to the intersection form}
Use the actual marked quotient \(G=G_3\) and
the arithmetic boundary map of D.54. Let
\(n=576\cdot3^{97}\). The subgroup \(K_3\)
is the fundamental group of the marked
oriented surface with
\[
 \begin{gathered}
 c=n/36\quad\hbox{cusps},\qquad
 2g=2+5n/36,\qquad
 \dim V=1+n/6,\\
 \mathcal R=\langle[g t^{36}g^{-1}]:
                 g\in G/\langle t\rangle\rangle
                 \subset V,\qquad
 \dim\mathcal R=c-1,\\
 H=V/\mathcal R,\qquad \dim H=2g .
 \end{gathered}
 \tag{D.55.1}
\]
Here \(V=H_1(K_3,\mathbb F_2)\), and the
projection \(\pi:V\to H\) is the actual
homology map of compactification.
The oriented surface determines its
nondegenerate intersection bivector
\(\Theta\in\bigwedge^2 H\), the inverse
of its alternating intersection form.

There is a specified linear map
\[
 \begin{aligned}
 \mathcal N:Z_G&\longrightarrow
                      (\bigwedge\nolimits^2H)^G,\\
 [z]&\longmapsto
       \bigwedge\nolimits^2\pi
                      \left(\sum_{g\in G}gz\right),
       \qquad Z=W\wedge V .
 \end{aligned}
 \tag{D.55.2}
\]
Every modulo-two pair parameter \(p\in\mathcal P\)
of D.54 gives the following equality of
actual maps:
\[
 \boxed{\quad
 \mathcal N(q_{\alpha_0}(p))
       =(\bigwedge\nolimits^2 T_{p,c})(\Theta)-\Theta .
       \quad}
 \tag{D.55.3}
\]
The map \(T_p=\overline T_0+D_p\) fixes
\(\mathcal R\) pointwise and induces
\(T_{p,c}\) on \(H\). Neither its
invertibility nor a zero of the arithmetic
boundary map is assumed in this equality.

The complete local image
\(\mathcal L=\ell(\mathcal K)\) of D.54
lies in \(\ker\mathcal N\).
Consequently the exact remaining comparison is
\[
 \begin{gathered}
 \overline{\mathcal N}:Z_G/\mathcal L
                \longrightarrow\operatorname{im}\mathcal N,\\
 0\longrightarrow
       \mathcal E_{\mathcal N}
       \longrightarrow Z_G/\mathcal L
       \xrightarrow{\ \overline{\mathcal N}\ }
                \operatorname{im}\mathcal N
       \longrightarrow0,\qquad
 \mathcal E_{\mathcal N}
       =\ker\mathcal N/\mathcal L .
 \end{gathered}
 \tag{D.55.4}
\]
No vanishing assertion for the actual
\(\mathcal E_{\mathcal N}\) is made.
Its two explicit sources are identified
by the exact sequence
\[
 \begin{gathered}
 \nu:Z_G\longrightarrow Z^G,\qquad
       [z]\longmapsto\sum_{g\in G}gz,\qquad
 p_C=(\bigwedge\nolimits^2\pi)|_Z,\\
 0\longrightarrow\widehat H^{-1}(G,Z)
       \longrightarrow\ker\mathcal N
       \xrightarrow{\ \nu\ }
          \operatorname{im}\nu\cap\ker p_C
       \longrightarrow0 .
 \end{gathered}
 \tag{D.55.5}
\]
Thus the norm kernel and the boundary
projection kernel remain in the
comparison, with their literal maps.

The actual zero fiber of
\(\overline q_{\alpha_0}\) consists
entirely of invertible homology actions.
Generation imposes no additional
restriction on that fiber: the compact
map preserves a nondegenerate bivector
and the cusp radical is already fixed.
The zero fiber itself remains
unevaluated, as does the full profinite
realization.

There is also a complete ordinary
relative-homology calculation for
the actual arithmetic base. With trivial
\(\mathbb F_2\) coefficients and the
marked peripheral collections
\[
 \mathcal S_G=
       (\langle s\rangle,\langle r\rangle,\langle t\rangle),
 \qquad
 \mathcal S_B=
       (\langle s\rangle,\langle r\rangle,
                       \overline{\langle t\rangle}),
 \tag{D.55.6}
\]
the cusp orders are \(36\) and
\(\widehat{\mathbb Z}\), respectively.
The exact maps are
\[
 \begin{gathered}
 H_2(Q,\mathbb Z)=C_2,\qquad
 Q^{\rm ab}=C_6,\qquad
 H_2(G,\mathbb F_2)\simeq\mathbb F_2^2,\\
 H_2(\langle s\rangle,\mathbb F_2)
    \oplus H_2(\langle t\rangle,\mathbb F_2)
       \xrightarrow{\ \sim\ }H_2(G,\mathbb F_2),\\
 H_2(\widehat B,\mathcal S_B;\mathbb F_2)
       \xrightarrow{\ \sim\ }
       H_2(G,\mathcal S_G;\mathbb F_2)
       \simeq\mathbb F_2 .
 \end{gathered}
 \tag{D.55.7}
\]
The two actual central double covers
\(\mathrm{SL}_2(\mathbb Z/12)\to Q\)
and D.48's \(\mathcal E\to Q\)
give the dual restriction matrix
\[
 \begin{array}{c|cc}
 &\langle s\rangle&\langle t\rangle\\ \hline
 \mathrm{SL}_2(\mathbb Z/12)&1&0\\
 \mathcal E&0&1
 \end{array}.
 \tag{D.55.8}
\]
Their pulled-back lift orders on \(G\)
are \((4,36)\) and \((2,72)\).
The relative orientation map is
therefore computed using the original
markings, rather than only abstract
dimensions.

For the separate order-\(1080\) family,
the norm in (D.55.2) has rank \(9\)
on the \(13\)-coordinate obstruction
space. Its kernel is exactly the
four-dimensional local image.
All \(18\) linear and all \(153\)
quadratic coefficients in (D.55.3)
are checked through that actual norm.
This proves the finite \(864\)-point
intersection comparison via the
specified morphism. It does not
assert equality of the arithmetic
norm kernel with its local image.
The manuscript passages are Part D(e),
section 51.16, folio 734 verso,
scan 174, and section 51.15,
folios 725--727, scans 156--160.
\end{statement}
\begin{proof}
\emph{The marked compactification and its homology map.}
The regular \(G\)-cover of the
oriented modular orbifold has
one edge for every element of \(G\),
with vertices indexed by
\(G/\langle s\rangle\) and
\(G/\langle r\rangle\).
The rotations are the actual
right multiplications by \(s,r\).
Their product gives the cusp cycles.
D.53 checked that \(t\) has exact
order \(36\) in \(G\); hence there
are \(c=n/36\) cusp cycles.
The torsion-free kernel has rank
\(\rho=1+n/6\). Capping these oriented
boundary components gives
\(2g=\rho-c+1=2+5n/36\).

The integral homology exact sequence
of this compactification identifies
the boundary sublattice with the
free group on the \(c\) oriented
cusp classes modulo their one sum.
It is a primitive sublattice.
Reduction modulo two proves the
dimension and quotient in (D.55.1);
in particular there is no additional
modulo-two boundary relation.
An oriented handle and boundary
presentation identifies the
intersection bivector as
\[
                  \Theta=\sum_{i=1}^{g}a_i\wedge b_i .
 \tag{D.55.9}
\]
It is independent of that presentation,
since it is the inverse of the
intersection form on compact homology.
All deck maps preserve the orientation,
so the quotient map and \(\Theta\)
are \(G\)-equivariant and \(G\)-invariant,
respectively.

\emph{The full cusp norm.}
Use the actual class-two lift of
the chosen pair parameter from D.54.
Write \(\varphi\) for its map from
the marked source to \(\mathcal T_2\).
It induces identity on \(G\), sends
the two torsion generators to their
specified conjugates, and has
\(\varphi(t)=zt\), \(z\in Z\).
These are actual endomorphisms of
the class-two target; no generation
assumption enters the argument.

Because \(Z\) is central in the
kernel \(D\), its conjugates by
powers of \(t\) commute. Consequently
\[
            (zt)^{36}
             =\left(\sum_{j=0}^{35}t^jz\right)t^{36}.
 \tag{D.55.10}
\]
Choose one marked coset representative
\(g_i\) for each cusp. The original
cusp loop is \(c_i=g_it^{36}g_i^{-1}\).
The element \(\varphi(g_i)g_i^{-1}\)
belongs to \(D\), since the base action
is identity. Thus the image of \(c_i\)
differs from its original conjugacy
class by the central term
\[
                         g_i\sum_{j=0}^{35}t^jz .
 \tag{D.55.11}
\]
The remaining conjugation correction
has exterior coordinate
\([u_i]\wedge[c_i]\) for a kernel
element \(u_i\). It maps to zero
under \(\bigwedge^2\pi\), because
the cusp class is in \(\mathcal R\).
In abelian homology the central error
also disappears. Hence every cusp
class is fixed, proving the asserted
pointwise action on \(\mathcal R\).

Sum (D.55.11) over the cusp cosets.
The representatives followed by the
thirty-six powers of \(t\) enumerate
each element of \(G\) once. The total
central error is therefore exactly
\(\sum_{g\in G}gz\). This proves the
literal full-group norm in (D.55.2).
The map is well-defined on
coinvariants since the sum is
unchanged by translating \(z\).
Its image is invariant after the
equivariant compact projection.

\emph{The orientation identity.}
Choose the integral punctured-surface
presentation with handles \(a_i,b_i\)
and boundary loops \(c_1,\ldots,c_c\):
\[
          \prod_{i=1}^{g}[a_i,b_i]\,
                    c_1\cdots c_c=1 .
 \tag{D.55.12}
\]
The free basis consists of the
handles and the first \(c-1\)
boundary generators. There is an
actual homomorphism of its free
class-two exponent-four quotient
which kills these \(c-1\) generators
and retains the handles. On central
coordinates it is exactly
\(\bigwedge^2\pi\).
The last boundary word then maps
to the inverse of the handle
commutator product. This construction
specifies the map used in the
following computation; it does not
assume that \(\varphi\) descends
through a compact surface relation.

Apply \(\varphi\) to (D.55.12) and
then use that homomorphism.
The handle commutators contribute
\((\bigwedge^2T_{p,c})(\Theta)\).
The original boundary words contribute
\(-\Theta\).
Their additional conjugation terms
vanish by the preceding paragraph,
and their central errors give
\(\mathcal N([z])\).
All surviving terms are central.
The resulting equality, in
characteristic two, is precisely
(D.55.3). Since \([z]\) is the
actual boundary map of D.54, this
also proves its compatibility with
all integral and local choices.

\emph{The exact residual and generation.}
A pair in \(\mathcal K\) has zero
homology restriction. Adding it
changes the boundary value by its
linear local image and leaves
\(T_{p,c}\) unchanged. Subtracting
the two orientation identities shows
\(\mathcal L\subset\ker\mathcal N\).
The kernel and image of the induced
linear map give (D.55.4).

For (D.55.5), the finite Tate group
in degree minus one is, by its
norm definition,
\[
 \widehat H^{-1}(G,Z)
   =\ker\!\left(\sum_{g\in G}g:Z\to Z\right)/I_GZ .
 \tag{D.55.13}
\]
It is exactly the kernel of \(\nu\).
The compact norm is the composition
\(p_C\nu\). Its kernel maps onto
\(\operatorname{im}\nu\cap\ker p_C\),
with the preceding Tate kernel.
This proves the claimed exact
sequence and retains both possible
kernel contributions.

On the zero fiber of the reduced
boundary map, (D.55.3) gives
\((\bigwedge^2T_{p,c})(\Theta)=\Theta\).
The support of the nondegenerate
bivector on the right is all of
\(H\). The support on the left is
contained in \(\operatorname{im}T_{p,c}\).
Thus \(T_{p,c}\) is onto, and hence
invertible in this finite-dimensional
space. Its full map \(T_p\) fixes
\(\mathcal R\) and induces this
isomorphism on the quotient; it is
therefore invertible on \(V\).
The complete homology lift is
invertible over \(\mathbb Z_2\),
and the finite class-two lift
generates by the Frattini criterion.
This proves the generation assertion
about the exact zero fiber without
asserting its nonemptiness.

\emph{The ordinary relative-homology computation.}
The discrete kernel \(N\) of
\(B=C_2*C_3\to Q\) has free rank \(97\).
D.26's integral lattice calculation,
reconstructed in
\texttt{d\_arithmetic\_\allowbreak{}relative\_h2.py},
gives
\[
                    (N^{\rm ab})_Q=C_2 .
 \tag{D.55.14}
\]
The exact Smith factors are
\((1^{96},2)\); appending the
actual cusp column \([t^{12}]\)
gives \((1^{97})\).
The marked abelianization map
\(Q\to C_6\) sends \(s\) to \(3\)
and \(r\) to \(2\). It is onto
and equals the abelianization,
because \(B^{\rm ab}=C_6\).
The homology five-term sequence,
with \(H_2(B,\mathbb Z)=0\),
therefore identifies \(H_2(Q,\mathbb Z)\)
with (D.55.14).

The universal coefficient sequence
has the two one-dimensional terms
\(H_2(Q,\mathbb Z)\otimes\mathbb F_2\)
and \(\operatorname{Tor}(C_6,\mathbb F_2)\).
Hence \(H_2(Q,\mathbb F_2)\)
has dimension two. The kernel of
\(G\to Q\) is the actual odd group
of order \(3^{97}\). Averaging its
bar resolution contracts all
positive-degree homology over
\(\mathbb F_2\). The induced map
on modulo-two homology is therefore
an isomorphism, giving the first
line of (D.55.7).

To compute the peripheral map,
use the two stated central double
covers. In \(\mathrm{SL}_2(\mathbb Z/12)\)
the marked lifts of \(s,r,t\)
have orders \(4,3,12\). In the
actual central cover \(\mathcal E\)
they have orders \(2,3,24\).
The program checks the full marked
state spaces and their projection
maps, including \(t^{12}\) as the
nontrivial central element of
\(\mathcal E\).
Pulling both extensions back to
\(G\) gives orders \(4,36\)
and \(2,72\) on \(s,t\).
A central double cover of an
even cyclic group has nonzero
restriction class exactly when
its generator lift has doubled
order. Thus the dual restriction
matrix is (D.55.8). It is invertible,
so the displayed map on \(H_2\)
is an isomorphism.

The \(r\)-subgroup has odd order,
so its positive modulo-two homology
vanishes. In the relative homology
sequence, the peripheral \(H_2\)
map is the just-proved isomorphism.
The next kernel is the one-dimensional
kernel of
\[
 H_1(\langle s\rangle,\mathbb F_2)
  \oplus H_1(\langle t\rangle,\mathbb F_2)
       \longrightarrow H_1(G,\mathbb F_2),
       \qquad(a,b)\longmapsto a+b .
 \tag{D.55.15}
\]
This proves the finite relative
dimension and specifies its boundary.
For \(\widehat B\), the free-product
homology sequence identifies its
\(H_2\) with that of \(C_2\).
The procyclic cusp has zero \(H_2\)
and one-dimensional \(H_1\).
The same relative sequence has the
same boundary kernel (D.55.15).
The original quotient preserves
both generators, so induces identity
on that kernel. This proves the
marked relative orientation
isomorphism in (D.55.7).

\emph{The exact separate finite norm certificate.}
The program
\texttt{d\_cm\_boundary\_norm.py}
reconstructs the order-\(1080\)
marked graph and the actual
intersection form of D.51.
It constructs the \(35\)-dimensional
cusp radical and the rank-\(146\)
compact quotient.
The chosen complement, projection
matrix and inverse-form bivector
are all retained. The code verifies
the actual \(s,r\) deck actions,
the compact projection and
orientation invariance.

It then reconstructs the exact
\(3600\)-dimensional invariant
exterior basis and the thirteen
original obstruction functionals.
Exact elimination supplies thirteen
representatives dual to those
functionals. The full group norm
is evaluated through literal
subgroup transversals of orders
\(3,4,5,3,2\), with the central
order-three average separately
checked. These are the actual
\(A_5\), \(S_3\), and central
\(C_3\) factors of this finite
source. They enumerate the full
group norm, not an unrelated
averaging operator.

After compact projection the norm
has rank nine. It kills the four
independent local-image vectors,
so its kernel is exactly their
span. Every one of the eighteen
linear coefficients and the
one hundred fifty-three quadratic
coefficients equals the corresponding
compact bivector defect through
this norm. This proves the stated
finite morphism and explains the
zero-locus equality of D.54.
The arithmetic maps above are
proved symbolically on their actual
source; the finite kernel equality
is not transferred to \(G_3\).
\end{proof}
